\begin{abstract}
Энэхүү судалгааны ажил нь жижиглэн худалдааны салбарт дэлгүүрийн лангуун дээрх бараа бүтээгдэхүүнийг зургаас автоматаар таньж, тооллого болон аудитын шийдвэр гаргалтыг хөнгөвчлөх системийг хөгжүүлэхэд оршино. Орчин үед дэлгүүрийн бараа материалын тооллого нь ихэвчлэн хүний оролцоотойгоор хийгдэж байгаа бөгөөд энэ нь цаг хугацаа их шаардсан, хөдөлмөрийн үр ашиггүй хэлбэрээс гадна алдаа гарах магадлал өндөр байдаг сул талтай. Бид энэхүү асуудлыг шийдвэрлэхийн тулд YOLO (You Only Look Once) архитектурын сүүлийн үеийн хувилбар болох YOLOv8 гүн сургалтын объектыг илрүүлэх загвар, FastAPI дээр суурилсан өндөр хурдны сүлжээний үйлчилгээ (Backend), болон хэрэглэгчид зориулсан олон платформт гар утасны Flutter аппликейшнээс бүрдсэн цогц мэдээллийн системийг санал болгож байна. Системийн ажиллах зарчмын хувьд гар утасны аппликейшнээр дамжуулан авсан лангууны зургийг сервер рүү илгээж, YOLOv8 загвар ашиглан зурган дээрх барааны байршил, төрлийг олж таних ба танигдсан үр дүнг өгөгдлийн санд урьдчилан тодорхойлсон хүлээгдэж буй барааны жагсаалттай харьцуулан автоматаар аудит хийх алгоритмыг хэрэгжүүлсэн. Туршилтын үр дүнгээс харахад сургагдсан загвар нь mAP50 хэмжүүрээр өндөр нарийвчлал үзүүлсэн ба нэг зургийг боловсруулах дундаж хугацаа нь үр дүнтэй (real-time) байлаа. Ингэснээр хөгжүүлсэн систем нь хүний оролцоог багасгаж, дэлгүүрийн менежментийн болон бараа нийлүүлэлтийн хяналт шалгалтын үр ашгийг эрс сайжруулах бүрэн боломжтойг батлан харуулж байна.
\end{abstract}
